\documentclass[11pt, a4paper]{article}

\usepackage[utf8]{inputenc}
\usepackage[T1]{fontenc}
\usepackage{geometry}
\geometry{margin=1in}
\usepackage{amsmath, amssymb, amsthm}
\usepackage{graphicx}
\usepackage{hyperref}
\usepackage{xcolor}
\usepackage{listings}
\usepackage{booktabs}
\usepackage{enumitem}

\definecolor{codegreen}{rgb}{0,0.6,0}
\definecolor{codegray}{rgb}{0.5,0.5,0.5}
\definecolor{codepurple}{rgb}{0.58,0,0.82}
\definecolor{backcolour}{rgb}{0.95,0.95,0.92}

\lstdefinestyle{mystyle}{
    backgroundcolor=\color{backcolour},   
    commentstyle=\color{codegreen},
    keywordstyle=\color{magenta},
    numberstyle=\tiny\color{codegray},
    stringstyle=\color{codepurple},
    basicstyle=\ttfamily\footnotesize,
    breakatwhitespace=false,         
    breaklines=true,                 
    captionpos=b,                    
    keepspaces=true,                 
    numbers=left,                    
    numbersep=5pt,                  
    showspaces=false,                
    showstringspaces=false,
    showtabs=false,                  
    tabsize=2,
    language=Python
}
\lstset{style=mystyle}

\title{\textbf{Lead-Lag Detection in Crypto Markets}\\[0.5em]\large From Spectral Methods to Temporal Graph Neural Networks}
\author{Project Specification}
\date{\today}

\begin{document}

\maketitle

\begin{abstract}
This project compares four approaches to discovering lead-lag relationships in cryptocurrency markets using tick-level trade data: (1) spectral analysis via the Hilbert transform, (2) Graph Attention Networks (GAT), (3) Temporal Graph Attention Networks (TGAT), and optionally (4) Sheaf Neural Networks. Unlike traditional equity markets, crypto trades 24/7 with publicly available tick data, making it ideal for studying high-frequency information propagation. We investigate whether BTC leads altcoins, how quickly price discovery propagates across assets, and whether learned attention patterns recover spectral phase relationships.
\end{abstract}

\tableofcontents
\newpage

%==============================================================================
\section{Motivation}
%==============================================================================

\subsection{Why Crypto?}

\begin{enumerate}
    \item \textbf{Free tick data}: Binance provides historical trade-level data going back years, no payment required
    \item \textbf{24/7 trading}: No overnight gaps, no market open/close effects
    \item \textbf{Natural event structure}: Each trade is a discrete event, fitting TGAT's design perfectly
    \item \textbf{Genuine lead-lag}: BTC is widely believed to lead altcoins; we can test this quantitatively
    \item \textbf{Fast dynamics}: Lead-lag effects occur over seconds to minutes, not days
\end{enumerate}

\subsection{The Core Question}

When BTC moves, how quickly do ETH, SOL, and other assets follow? Is this propagation structure:
\begin{itemize}
    \item Recoverable from frequency-domain analysis (spectral)?
    \item Learnable by a contemporaneous attention model (GAT)?
    \item Better captured by temporal attention (TGAT)?
    \item Expressible as phase shifts in a sheaf structure?
\end{itemize}

\subsection{What We'll Compare}

\begin{center}
\begin{tabular}{lllll}
\toprule
\textbf{Method} & \textbf{Temporal?} & \textbf{Learned?} & \textbf{Output} & \textbf{Fit to Tick Data} \\
\midrule
Spectral (Hilbert) & Yes (frequency domain) & No & Phases $\phi_{ij}$ & Good (after resampling) \\
GAT & No (contemporaneous) & Yes & Attention $\alpha_{ij}$ & Poor (snapshot-based) \\
GAT + lags & Implicit & Yes & Attention $\alpha_{ij}$ & Okay (hack) \\
TGAT & Yes (explicit) & Yes & Temporal attention & Native fit \\
Sheaf NN & Yes (by design) & Yes & Learned phases & Good (after resampling) \\
\bottomrule
\end{tabular}
\end{center}

%==============================================================================
\section{Data}
%==============================================================================

\subsection{Asset Universe}

Ten liquid trading pairs on Binance (all against USDT):

\begin{center}
\begin{tabular}{llll}
\toprule
\textbf{Ticker} & \textbf{Name} & \textbf{Category} & \textbf{Expected Role} \\
\midrule
BTC & Bitcoin & Store of value & Leader \\
ETH & Ethereum & Smart contracts & Fast follower \\
BNB & Binance Coin & Exchange token & Exchange-specific \\
SOL & Solana & Alt L1 & Follower \\
XRP & Ripple & Payments & Idiosyncratic \\
ADA & Cardano & Alt L1 & Follower \\
DOGE & Dogecoin & Meme & Retail-driven \\
MATIC & Polygon & L2 & ETH-correlated \\
LINK & Chainlink & Oracle & DeFi-correlated \\
AVAX & Avalanche & Alt L1 & Follower \\
\bottomrule
\end{tabular}
\end{center}

\subsection{Data Source}

Binance historical trade data: \url{https://data.binance.vision/}

\begin{lstlisting}[language=bash]
# Download March 2024 trades for BTCUSDT
wget https://data.binance.vision/data/spot/daily/trades/BTCUSDT/BTCUSDT-trades-2024-03-01.zip
\end{lstlisting}

Each file contains:
\begin{center}
\begin{tabular}{ll}
\toprule
\textbf{Field} & \textbf{Description} \\
\midrule
\texttt{trade\_id} & Unique trade identifier \\
\texttt{price} & Execution price \\
\texttt{qty} & Quantity traded \\
\texttt{quote\_qty} & Quote asset quantity \\
\texttt{time} & Unix timestamp (milliseconds) \\
\texttt{is\_buyer\_maker} & Trade direction \\
\bottomrule
\end{tabular}
\end{center}

\subsection{Data Volume}

Typical volumes:
\begin{itemize}
    \item BTC/USDT: $\sim$1--2 million trades/day
    \item ETH/USDT: $\sim$500k--1M trades/day
    \item Smaller altcoins: $\sim$100k--500k trades/day
\end{itemize}

For one week of data across 10 assets: $\sim$50--100 million trades. We'll subsample.

\subsection{Preprocessing}

\begin{enumerate}
    \item \textbf{Download}: Fetch daily ZIP files for each asset
    \item \textbf{Parse}: Load into DataFrames, convert timestamps
    \item \textbf{Aggregate}: Create uniform time bars (e.g., 1-second or 5-second bars)
    \item \textbf{Align}: Ensure all assets have same timestamp index
    \item \textbf{Compute returns}: Log returns for each bar
\end{enumerate}

\begin{lstlisting}[language=Python]
import pandas as pd
import numpy as np

def load_trades(filepath):
    """Load Binance trade data from ZIP."""
    df = pd.read_csv(
        filepath,
        compression='zip',
        names=['trade_id', 'price', 'qty', 'quote_qty', 
               'time', 'is_buyer_maker', 'is_best_match']
    )
    df['time'] = pd.to_datetime(df['time'], unit='ms')
    df = df.set_index('time')
    return df

def aggregate_to_bars(trades, freq='5S'):
    """Aggregate trades to OHLCV bars."""
    bars = trades['price'].resample(freq).ohlc()
    bars['volume'] = trades['qty'].resample(freq).sum()
    bars['num_trades'] = trades['price'].resample(freq).count()
    bars = bars.dropna()
    return bars

def compute_returns(bars):
    """Compute log returns from close prices."""
    return np.log(bars['close'] / bars['close'].shift(1)).dropna()
\end{lstlisting}

\subsection{Time Period}

\begin{itemize}
    \item \textbf{Development}: 1 week of data (March 1--7, 2024)
    \item \textbf{Full analysis}: 1 month if compute permits
    \item \textbf{Train/Val/Test}: 60\%/20\%/20\% chronological split
\end{itemize}

%==============================================================================
\section{Method 1: Spectral Lead-Lag}
%==============================================================================

\subsection{Overview}

Same approach as before, but on higher-frequency data:
\begin{enumerate}
    \item Aggregate trades to fixed bars (5-second)
    \item Compute log returns
    \item Apply Hilbert transform to get analytic signals
    \item Compute complex correlation matrix
    \item Extract global phases via Angular Synchronization
\end{enumerate}

\subsection{Resolution Choice}

\begin{center}
\begin{tabular}{lll}
\toprule
\textbf{Bar Size} & \textbf{Bars/Day} & \textbf{Lead-Lag Horizon} \\
\midrule
1 second & 86,400 & Sub-second to seconds \\
5 seconds & 17,280 & Seconds to minutes \\
30 seconds & 2,880 & Minutes \\
1 minute & 1,440 & Minutes to hours \\
\bottomrule
\end{tabular}
\end{center}

Recommendation: Start with \textbf{5-second bars}. Fine enough to capture fast dynamics, coarse enough to reduce noise.

\subsection{Implementation}

\begin{lstlisting}[language=Python]
from scipy import signal

def compute_spectral_phases(returns_df):
    """
    Compute lead-lag phases via Hilbert transform.
    
    returns_df: (T, N) DataFrame of log returns, columns are assets
    Returns: (N,) array of global phases
    """
    # Analytic signals
    analytic = pd.DataFrame(
        signal.hilbert(returns_df.values, axis=0),
        index=returns_df.index,
        columns=returns_df.columns
    )
    
    # Standardize
    Z = (analytic - analytic.mean()) / analytic.std()
    
    # Complex correlation matrix
    C = Z.values.conj().T @ Z.values / len(Z)
    
    # Principal eigenvector
    eigenvalues, eigenvectors = np.linalg.eigh(C)
    v1 = eigenvectors[:, -1]
    
    # Extract and center phases
    theta = np.angle(v1)
    theta -= np.angle(np.mean(np.exp(1j * theta)))
    
    return theta, C
\end{lstlisting}

\subsection{Expected Output}

A ranking like:
\begin{center}
\begin{tabular}{cl}
\toprule
\textbf{Rank} & \textbf{Asset} \\
\midrule
1 (leads) & BTC \\
2 & ETH \\
3 & BNB \\
4 & SOL \\
$\vdots$ & $\vdots$ \\
10 (lags) & DOGE \\
\bottomrule
\end{tabular}
\end{center}

If BTC doesn't come out on top, that's interesting---maybe ETH leads in DeFi-correlated moves?

%==============================================================================
\section{Method 2: GAT (Baseline)}
%==============================================================================

\subsection{Why Include It?}

GAT isn't designed for temporal data, but it's a useful baseline:
\begin{itemize}
    \item If GAT works well, maybe temporal structure isn't crucial
    \item If GAT fails, we have evidence that proper temporal modeling matters
    \item It's simple to implement and understand
\end{itemize}

\subsection{Setup}

\textbf{Graph}: Fully connected (10 nodes, 90 directed edges)

\textbf{Node features} at each time $t$ (per bar):
\begin{itemize}
    \item \texttt{ret\_1}: Return this bar
    \item \texttt{ret\_5}: Cumulative return over last 5 bars
    \item \texttt{ret\_20}: Cumulative return over last 20 bars
    \item \texttt{vol\_20}: Rolling volatility (20 bars)
    \item \texttt{volume\_zscore}: Volume relative to recent average
    \item \texttt{num\_trades}: Number of trades this bar
\end{itemize}

\textbf{With lagged features} (Option B from before):
\begin{itemize}
    \item \texttt{ret\_lag1}: Return of previous bar
    \item \texttt{ret\_lag2}: Return two bars ago
    \item \texttt{neighbor\_ret\_lag1\_mean}: Average neighbor return last bar
\end{itemize}

\textbf{Target}: Next-bar return sign (binary classification)

\subsection{Architecture}

\begin{lstlisting}[language=Python]
import torch
import torch.nn.functional as F
from torch_geometric.nn import GATConv

class CryptoGAT(torch.nn.Module):
    def __init__(self, in_channels, hidden_channels=32, heads=4):
        super().__init__()
        self.conv1 = GATConv(in_channels, hidden_channels, heads=heads, concat=True)
        self.conv2 = GATConv(hidden_channels * heads, hidden_channels, heads=1, concat=False)
        self.lin = torch.nn.Linear(hidden_channels, 1)
    
    def forward(self, x, edge_index, return_attention=False):
        x, (edge_index_1, alpha_1) = self.conv1(
            x, edge_index, return_attention_weights=True
        )
        x = F.elu(x)
        x = F.dropout(x, p=0.3, training=self.training)
        
        x, (edge_index_2, alpha_2) = self.conv2(
            x, edge_index, return_attention_weights=True
        )
        x = F.elu(x)
        
        out = self.lin(x)
        
        if return_attention:
            return out, alpha_1
        return out
\end{lstlisting}

\subsection{Attention Extraction}

After training, extract "who attends to whom":
\begin{lstlisting}[language=Python]
def get_attention_ranking(model, data):
    """Compute attention inflow for each node."""
    model.eval()
    with torch.no_grad():
        _, alpha = model(data.x, data.edge_index, return_attention=True)
    
    # alpha: (num_edges, num_heads) -> average over heads
    alpha_mean = alpha.mean(dim=1)
    
    # Sum attention flowing INTO each node
    num_nodes = data.x.shape[0]
    inflow = torch.zeros(num_nodes)
    for i, (src, tgt) in enumerate(data.edge_index.T):
        inflow[tgt] += alpha_mean[i]
    
    return inflow.numpy()
\end{lstlisting}

%==============================================================================
\section{Method 3: TGAT}
%==============================================================================

\subsection{Why TGAT Fits}

TGAT was designed for exactly this: irregular events with timestamps. Each trade is an event, and TGAT can attend over historical events for each node.

\subsection{Data Format}

TGAT expects a stream of events, not snapshots:

\begin{center}
\begin{tabular}{llll}
\toprule
\textbf{Source} & \textbf{Target} & \textbf{Time} & \textbf{Features} \\
\midrule
BTC & BTC & 1709312400.123 & [price, volume, ...] \\
ETH & ETH & 1709312400.456 & [price, volume, ...] \\
BTC & BTC & 1709312401.001 & [price, volume, ...] \\
\bottomrule
\end{tabular}
\end{center}

Self-edges represent "node $i$ updated at time $t$." Cross-edges can represent "price correlation links."

\subsection{Two Formulations}

\textbf{Option A: Trade events only}
\begin{itemize}
    \item Each trade is an event on that asset's node
    \item Cross-asset edges connect correlated assets
    \item TGAT attends over each asset's recent trade history
\end{itemize}

\textbf{Option B: Regular sampling with event structure}
\begin{itemize}
    \item Sample at fixed intervals (5 seconds)
    \item Treat each sample as an "event" for all nodes
    \item Simpler, but loses some temporal granularity
\end{itemize}

Recommendation: Start with Option B (simpler), then try Option A if time permits.

\subsection{Architecture}

\begin{lstlisting}[language=Python]
# Using PyTorch Geometric Temporal
from torch_geometric_temporal.nn.attention import TGAT

class CryptoTGAT(torch.nn.Module):
    def __init__(self, in_channels, hidden_channels=32, num_heads=4):
        super().__init__()
        self.tgat = TGAT(
            in_channels=in_channels,
            out_channels=hidden_channels,
            heads=num_heads,
        )
        self.lin = torch.nn.Linear(hidden_channels * num_heads, 1)
    
    def forward(self, x, edge_index, edge_time):
        """
        x: (num_nodes, in_channels) - node features
        edge_index: (2, num_edges) - graph structure
        edge_time: (num_edges,) - timestamp of each edge
        """
        h = self.tgat(x, edge_index, edge_time)
        return self.lin(h)
\end{lstlisting}

\subsection{Temporal Attention Interpretation}

TGAT produces attention weights $\alpha_{ij}^{(t, \tau)}$: "how much does node $i$ at time $t$ attend to node $j$ at time $\tau < t$?"

For lead-lag:
\begin{itemize}
    \item High $\alpha_{\text{ETH} \leftarrow \text{BTC}}^{(t, t-1)}$: BTC's recent past matters for ETH now
    \item Compare $\alpha_{i \leftarrow j}$ vs $\alpha_{j \leftarrow i}$ to determine who leads
\end{itemize}

%==============================================================================
\section{Method 4: Sheaf Neural Networks (Optional)}
%==============================================================================

\subsection{The Idea}

Model lead-lag as phase shifts in a $U(1)$ sheaf:
\begin{itemize}
    \item Each node has a complex-valued state
    \item Restriction maps are phase rotations $e^{i\phi_{ij}}$
    \item Learn the phases end-to-end
\end{itemize}

\subsection{Why It's Hard}

\begin{enumerate}
    \item Complex tensor support in PyTorch is incomplete
    \item Phases live on a circle, not $\mathbb{R}$---gradient descent doesn't respect this
    \item Antisymmetry constraint ($\phi_{ij} = -\phi_{ji}$) must be enforced
    \item Initialization matters: random phases $\to$ no learning
    \item Degenerate solutions (all phases $\to$ 0) require careful regularization
\end{enumerate}

\subsection{Implementation Sketch}

\begin{lstlisting}[language=Python]
class SheafConv(torch.nn.Module):
    def __init__(self, num_nodes, in_channels, out_channels):
        super().__init__()
        # Learnable phases (upper triangle only, for antisymmetry)
        n_pairs = num_nodes * (num_nodes - 1) // 2
        self.phase_params = torch.nn.Parameter(torch.zeros(n_pairs))
        self.lin = torch.nn.Linear(in_channels, out_channels)
    
    def get_phase_matrix(self, num_nodes):
        """Build antisymmetric phase matrix."""
        idx = torch.triu_indices(num_nodes, num_nodes, offset=1)
        phases = torch.zeros(num_nodes, num_nodes)
        phases[idx[0], idx[1]] = self.phase_params
        phases[idx[1], idx[0]] = -self.phase_params
        return phases
    
    def forward(self, x, adj):
        num_nodes = x.shape[0]
        x = self.lin(x)
        
        # Magnetic adjacency
        phases = self.get_phase_matrix(num_nodes)
        magnetic_adj = adj * torch.exp(1j * phases)
        
        # Normalized diffusion
        deg = adj.sum(dim=1)
        deg_inv_sqrt = torch.diag(1.0 / torch.sqrt(deg + 1e-8))
        norm_adj = deg_inv_sqrt @ magnetic_adj @ deg_inv_sqrt
        
        # Aggregate (complex-valued)
        x_complex = x.to(torch.cfloat)
        out = norm_adj @ x_complex
        
        return out.real
\end{lstlisting}

\subsection{When to Attempt}

Only after completing spectral, GAT, and TGAT. Budget 1--2 weeks for debugging.

%==============================================================================
\section{Comparison Analysis}
%==============================================================================

\subsection{Metric 1: Rank Correlation}

Compare lead-lag rankings across methods:

\begin{lstlisting}[language=Python]
from scipy.stats import spearmanr

# Rankings from each method
spectral_rank = np.argsort(-spectral_phases)  # Higher phase = leads
gat_rank = np.argsort(-gat_attention_inflow)
tgat_rank = np.argsort(-tgat_attention_inflow)

# Pairwise Spearman correlations
rho_spectral_gat, _ = spearmanr(spectral_rank, gat_rank)
rho_spectral_tgat, _ = spearmanr(spectral_rank, tgat_rank)
rho_gat_tgat, _ = spearmanr(gat_rank, tgat_rank)
\end{lstlisting}

\subsection{Metric 2: Prediction Performance}

\begin{center}
\begin{tabular}{lll}
\toprule
\textbf{Model} & \textbf{Task} & \textbf{Metrics} \\
\midrule
Baseline (logistic) & Predict next-bar sign & Accuracy, AUC \\
GAT & Predict next-bar sign & Accuracy, AUC \\
GAT + lags & Predict next-bar sign & Accuracy, AUC \\
TGAT & Predict next-bar sign & Accuracy, AUC \\
\bottomrule
\end{tabular}
\end{center}

\subsection{Metric 3: Qualitative Analysis}

\begin{enumerate}
    \item \textbf{Market Clock}: Polar plot of spectral phases
    \item \textbf{Attention heatmaps}: GAT and TGAT attention matrices
    \item \textbf{Lead-lag consistency}: Does BTC consistently lead across methods?
    \item \textbf{Regime dependence}: Do rankings change during high volatility?
\end{enumerate}

\subsection{Hypotheses}

Before running experiments:

\begin{enumerate}
    \item BTC leads all other assets (spectral phase highest)
    \item ETH is second (fast follower)
    \item DOGE is most lagging (retail-driven, slow to react)
    \item TGAT outperforms GAT on prediction (temporal structure matters)
    \item Spectral and TGAT rankings correlate ($\rho > 0.6$)
    \item GAT rankings correlate less with spectral ($\rho \approx 0.3$--$0.5$)
\end{enumerate}

%==============================================================================
\section{Implementation Plan}
%==============================================================================

\subsection{Week 1: Data Pipeline + Spectral}

\begin{center}
\begin{tabular}{lp{10cm}}
\toprule
\textbf{Day} & \textbf{Task} \\
\midrule
1 & Write download script for Binance historical data \\
2 & Load and parse trade data for 10 assets, one day \\
3 & Aggregate to 5-second bars, align timestamps \\
4 & Implement Hilbert transform pipeline, compute phases \\
5 & Create Market Clock visualization, validate BTC leads \\
\bottomrule
\end{tabular}
\end{center}

\textbf{Deliverables}:
\begin{itemize}
    \item \texttt{data/trades/*.parquet} --- raw trade data
    \item \texttt{data/bars\_5s.parquet} --- aggregated bars
    \item \texttt{results/spectral\_phases.csv}
    \item \texttt{figures/market\_clock.png}
\end{itemize}

\subsection{Week 2: GAT}

\begin{center}
\begin{tabular}{lp{10cm}}
\toprule
\textbf{Day} & \textbf{Task} \\
\midrule
1 & Create PyG Data objects from bar data \\
2 & Implement GAT model, verify forward pass \\
3 & Add lagged features (Option B) \\
4 & Train GAT, extract attention weights \\
5 & Compare attention ranking to spectral ranking \\
\bottomrule
\end{tabular}
\end{center}

\textbf{Deliverables}:
\begin{itemize}
    \item \texttt{models/gat.pt}
    \item \texttt{results/gat\_attention.csv}
    \item \texttt{figures/attention\_heatmap.png}
\end{itemize}

\subsection{Week 3: TGAT}

\begin{center}
\begin{tabular}{lp{10cm}}
\toprule
\textbf{Day} & \textbf{Task} \\
\midrule
1 & Install \texttt{torch\_geometric\_temporal}, understand API \\
2 & Reformat data for TGAT (event stream or batched snapshots) \\
3 & Implement TGAT model \\
4 & Train TGAT, extract temporal attention \\
5 & Compare TGAT ranking to spectral and GAT \\
\bottomrule
\end{tabular}
\end{center}

\textbf{Deliverables}:
\begin{itemize}
    \item \texttt{models/tgat.pt}
    \item \texttt{results/tgat\_attention.csv}
    \item \texttt{figures/temporal\_attention.png}
\end{itemize}

\subsection{Week 4: Comparison + Write-up}

\begin{center}
\begin{tabular}{lp{10cm}}
\toprule
\textbf{Day} & \textbf{Task} \\
\midrule
1 & Compute all comparison metrics (rank correlations, prediction accuracy) \\
2 & Create comparison visualizations \\
3 & Analyze disagreements between methods \\
4 & Write final report \\
5 & Clean code, documentation \\
\bottomrule
\end{tabular}
\end{center}

\textbf{Deliverables}:
\begin{itemize}
    \item \texttt{results/comparison\_metrics.json}
    \item \texttt{figures/ranking\_comparison.png}
    \item \texttt{report/final\_report.pdf}
\end{itemize}

\subsection{Week 5+ (Optional): Sheaf NN}

\begin{center}
\begin{tabular}{lp{10cm}}
\toprule
\textbf{Day} & \textbf{Task} \\
\midrule
1--2 & Implement SheafConv with learnable phases \\
3 & Handle complex gradients, phase projection \\
4 & Initialize from spectral phases, train \\
5 & Compare learned phases to spectral ground truth \\
6--7 & Debug, regularize, iterate \\
\bottomrule
\end{tabular}
\end{center}

%==============================================================================
\section{Technical Stack}
%==============================================================================

\begin{center}
\begin{tabular}{ll}
\toprule
\textbf{Component} & \textbf{Library} \\
\midrule
Data download & \texttt{requests}, \texttt{wget} \\
Data processing & \texttt{pandas}, \texttt{numpy}, \texttt{pyarrow} \\
Signal processing & \texttt{scipy.signal} \\
Graph neural networks & \texttt{torch\_geometric} \\
Temporal GNNs & \texttt{torch\_geometric\_temporal} \\
Visualization & \texttt{matplotlib}, \texttt{seaborn}, \texttt{networkx} \\
\bottomrule
\end{tabular}
\end{center}

\subsection{Installation}

\begin{lstlisting}[language=bash]
pip install pandas numpy scipy matplotlib seaborn networkx pyarrow
pip install torch torch_geometric
pip install torch_geometric_temporal
\end{lstlisting}

%==============================================================================
\section{Expected Results}
%==============================================================================

\subsection{Spectral Ranking (Hypothesized)}

\begin{center}
\begin{tabular}{cll}
\toprule
\textbf{Rank} & \textbf{Asset} & \textbf{Rationale} \\
\midrule
1 & BTC & Market leader, most liquidity \\
2 & ETH & Second largest, DeFi hub \\
3 & BNB & Exchange token, Binance-specific info \\
4 & SOL & High-beta alt L1 \\
5 & LINK & DeFi infrastructure \\
6 & AVAX & Alt L1 \\
7 & MATIC & ETH L2, follows ETH \\
8 & ADA & Slower community \\
9 & XRP & Idiosyncratic (legal issues) \\
10 & DOGE & Retail meme coin, slowest \\
\bottomrule
\end{tabular}
\end{center}

\subsection{Method Agreement (Hypothesized)}

\begin{center}
\begin{tabular}{lccc}
\toprule
& \textbf{Spectral} & \textbf{GAT} & \textbf{TGAT} \\
\midrule
Spectral & 1.0 & 0.4 & 0.7 \\
GAT & & 1.0 & 0.5 \\
TGAT & & & 1.0 \\
\bottomrule
\end{tabular}
\end{center}

TGAT should agree more with spectral (both capture temporal structure) while GAT captures something different (contemporaneous importance).

\subsection{Prediction Performance (Hypothesized)}

\begin{center}
\begin{tabular}{lcc}
\toprule
\textbf{Model} & \textbf{Accuracy} & \textbf{AUC} \\
\midrule
Baseline & 0.50 & 0.50 \\
GAT & 0.51 & 0.52 \\
GAT + lags & 0.52 & 0.54 \\
TGAT & 0.53 & 0.56 \\
\bottomrule
\end{tabular}
\end{center}

Modest improvements---crypto is noisy and efficient. But TGAT should show consistent edge.

%==============================================================================
\section{Success Criteria}
%==============================================================================

The project is successful if:

\begin{enumerate}
    \item \textbf{Data pipeline works}: Can download, parse, and aggregate Binance tick data
    \item \textbf{Spectral method produces sensible ranking}: BTC leads or there's a good explanation why not
    \item \textbf{GAT and TGAT train without diverging}: Produce non-trivial attention patterns
    \item \textbf{Comparison is quantified}: Rank correlations computed and interpreted
    \item \textbf{You understand all four methods}: Can explain Hilbert transform, GAT attention, TGAT temporal attention, and sheaf phases
\end{enumerate}

The project is \emph{not} judged by:
\begin{itemize}
    \item High prediction accuracy (crypto is efficient)
    \item Perfect method agreement (disagreement is interesting)
    \item Completing Sheaf NN (it's optional)
\end{itemize}

%==============================================================================
\section{References}
%==============================================================================

\begin{thebibliography}{99}

\bibitem{tgat} Xu, D., et al. ``Inductive Representation Learning on Temporal Graphs.'' ICLR, 2020.

\bibitem{gat} Veli\v{c}kovi\'{c}, P., et al. ``Graph Attention Networks.'' ICLR, 2018.

\bibitem{sheaf} Bodnar, C., et al. ``Neural Sheaf Diffusion.'' NeurIPS, 2022.

\bibitem{hilbert} Marple, S.L. ``Computing the Discrete-Time Analytic Signal via FFT.'' IEEE Trans.\ Signal Processing, 1999.

\bibitem{crypto_leadlag} Aslanidis, N., et al. ``Lead-Lag Relationships in Cryptocurrency Markets.'' Finance Research Letters, 2022.

\end{thebibliography}

\end{document}
